% TOP_PROBLEM: {"id": 6600009, "title": "U. Grimm: Diffraction of a Pinwheel Tiling — Problem", "category_id": 6, "category": {"id": 6, "name": "geometry", "display_name": "Geometry", "description": "Euclidean and non-Euclidean geometry, geometric structures.", "slug": "geometry", "order_index": 6, "created_at": "2026-07-31T15:26:25.671Z"}, "status": "open"}
% FIRST_POSTED: null
% ENTRY_KIND: statement_only
% Imported from: batch13.tex; UnsolvedMath ID 6600009
\documentclass[11pt]{article}
\usepackage[margin=1in]{geometry}
\usepackage{amsmath,amssymb,mathtools}
\usepackage{fontspec,unicode-math}
\setmainfont{Latin Modern Roman}
\setsansfont{Latin Modern Sans}
\setmathfont{Latin Modern Math}
\usepackage{xurl,hyperref}
\hypersetup{colorlinks=true,urlcolor=blue,linkcolor=blue}
\newcommand{\sref}[2]{\hyperlink{source:#1:#2}{[S#2]}}
\newcommand{\cataloguescope}[1]{\paragraph{Recorded mathematical question.}#1}
\begin{document}
% UnsolvedMath ID 6600009; source catalogue code AMR-065-0009
\setcounter{section}{6600008}
\section{U. Grimm: Diffraction of a Pinwheel Tiling — Problem}\label{6600009}
{\small\sffamily UnsolvedMath ID 6600009 · Geometry}
\subsection{Definitions and mathematical statement}

\paragraph{Shared notation.}$\mathbb N=\{1,2,\ldots\}$, and $\mathbb Z,\mathbb Q,\mathbb R,\mathbb C$ denote the integers, rationals, reals and complex numbers. Any other conventions are specified locally.


A pinwheel tiling is the Conway--Radin substitution tiling by right triangles of side lengths $1,2,\sqrt5$: inflating by $\sqrt5$ yields five congruent original triangles in the standard pinwheel subdivision. Choose on each triangle the point whose coordinates from the right-angle vertex along its perpendicular unit axes are $(1/2,1/2)$, and let $\Lambda$ be all these points. The uniform Dirac comb is $\omega=\sum_{x\in\Lambda}\delta_x$. Its volume-averaged autocorrelation is
\[\gamma=\lim_{R\to\infty}\frac{\omega|_{B_R}*\widetilde{\omega|_{B_R}}}{\operatorname{area}(B_R)},\]
in the vague topology; the tilde reflects and conjugates a measure. The positive diffraction measure is its Fourier transform $\widehat\gamma$, with kernel $e^{-2\pi i k\cdot x}$. This measure is rotationally invariant. An absolutely continuous diffraction component means a nonzero term $f(k)\,dk$ in the Lebesgue decomposition of $\widehat\gamma$ with respect to two-dimensional Lebesgue measure. It is distinct from a circular-ring component, which is singular as a measure on the plane.
\paragraph{Statement recorded in the supplied source.}
Does pinwheel diffraction have a nonzero absolutely continuous component? \sref{6600009}{2}
\subsection{Short English statement}
Is some pinwheel diffraction intensity spread with a density over the plane?
\subsection{Sources}
\begin{description}
\item[\hypertarget{source:6600009:1}{S1}] ULAM.AI and the UnsolvedMath Contributors, \emph{UnsolvedMath: A Curated Collection of Open Mathematics Problems}, record 6600009 (AMR-065-0009). CC BY 4.0. \url{https://huggingface.co/datasets/ulamai/UnsolvedMath}.
\item[\hypertarget{source:6600009:2}{S2}] Faustin Adiceam (compiler), \emph{Open Problems and Conjectures related to the Theory of Mathematical Quasicrystals}, arXiv:1604.06280v2 (2016), Section 2.3, construction and Problem 2.3.2, pp. 5--6. \url{https://arxiv.org/pdf/1604.06280}.
\end{description}
\label{end:6600009}
\end{document}
